\section{Dinámica de resistencia entre capas}

El análisis de los datos revela una dinámica de \textbf{resistencia bidireccional} que impide tanto la integración ascendente indiscriminada como la disolución descendente de estructuras ya consolidadas.

\subsection{«Resistencia ascendente» (o resistencia de abajo hacia arriba)}

Las capas inferiores (presentes locales intensificados) pueden resistir la reorganización estructural manteniendo \textbf{alta antisincronización}. Al sostener escalas de persistencia autónomas e independientes, los módulos generan fragmentación que actúa como barrera: una intensificación local no encuentra los canales relacionales necesarios para propagarse. La fragmentación funciona, por tanto, como un mecanismo de \textit{protección de la diversidad local} —a este fenómeno lo denominamos «resistencia ascendente»— frente a presiones globales de unificación.

\subsection{«Resistencia descendente» (o resistencia de arriba hacia abajo)}

Una vez alcanzada cierta coherencia entre los presentes locales (Capa 2 consolidada), el presente de nivel superior adquiere una \textbf{inercia ontológica}. Esta inercia le permite filtrar o absorber intensificaciones locales que emergen de forma todavía fragmentada. El sistema de mayor nivel no ignora completamente lo que ocurre en los módulos, pero puede resistir disrupciones que no logran reducir suficientemente la fragmentación. Esta resistencia —que proponemos llamar «resistencia descendente»— protege la estructura global ya alcanzada frente a fluctuaciones locales efímeras.

\subsection{Implicación ontológica}

La existencia de resistencias en ambas direcciones sugiere que el presente, una vez consolidado en capas superiores, goza de una \textit{autonomía relativa} respecto de los presentes locales. No se trata de un epifenómeno ni de una mera suma, sino de una capa con propiedades emergentes de integración y resistencia que retroactúan sobre los niveles inferiores. Esta autonomía relativa es compatible con la idea de que las transiciones entre capas constituyen los eventos discretos ontológicamente relevantes del tiempo extramental.

La antisincronización aparece así como el \textbf{operador de fragmentación} fundamental: su aumento incrementa la «resistencia ascendente»; su disminución reduce la «resistencia descendente» y abre la posibilidad de integración estructural.